\appendix

\section{Ontology Derivation Prompt}
\label{app:ontology_prompt}

The following prompt was used to query GPT-5.2, Gemini 3.1, and DeepSeek-V3.2
for candidate Wikidata QIDs and PIDs relevant to supply chain relation
extraction:

\begin{quote}
{\itshape 
Act as a Knowledge Graph Engineer. I am building a Relation Extraction dataset based on SCOR-DS and Enterprise Ontology.
I have the following seed Q-codes:
Actors: business (Q4830453), enterprise (Q688151), corporation (Q167037), vendor (Q11270215), third-party source (Q11832159), subsidiary (Q355209), government agency (Q327333), customer (Q22855), competitor (Q476635), retailer (Q200552).
Products and Materials:  product (Q2424752), goods (Q28877), manufactured product (Q1836700), raw material (Q192355), intermediate good (Q3955017)

Task: Identify 20-30 additional Wikidata Q-codes that represent:

    \begin{itemize}
        \item Functional Roles: (e.g., 'Wholesaler', 'Freight Forwarder', 'OEM', 'Contract Manufacturer', '3PL').
        \item Products \& Materials: (e.g., 'Intermediate good', 'Raw material', 'Electronic component', 'Finished good', 'Commodity').
        \item For each code, provide its Superclass QID to ensure it falls under either the 'Organization/Economic Entity' hierarchy or the 'Product/Physical Object' hierarchy.
    \end{itemize}

Organization/Constraint:

    \begin{itemize}
        \item Group the results into my 5 meta-categories: Core Enterprise, Supply Chain Partner, Organizational Unit, Market Actor, and Product \& Material.
        \item Exclude generic abstract classes like 'Entity' or 'Object'.
        \item Priority: Only provide classes that are high-level enough to have significant instances with Wikipedia pages.
    \end{itemize}}
\end{quote}


\section{Ontology Schema}
\label{app:ontology}

Table~\ref{tab:pid_mapping} lists the fifteen relation types in the
SC-RED ontology alongside their corresponding Wikidata property
identifiers and SCOR-DS Phases. Table~\ref{tab:extra_pids} lists the additional properties proposed alongside the corresponding SCOR-DS Phases.

\begin{table}[h]
\centering
\caption{Entity Types}
\label{tab:entity_types}
\begin{tabular}{ll}
\textbf{Entity Type} & \textbf{Description}\\
CoreEnterprise      & profit-seeking companies
                         (e.g. logistics company, freight forwarder, 
                          public company, defense contractor) \\
SupplyChainPartner  & external suppliers and service providers
                         (e.g. vendor, distributor, OEM, 
                          contract manufacturer) \\
OrganizationalUnit  & internal divisions and regulatory bodies
                         (e.g. government agency, regulatory body, 
                          joint venture) \\
MarketActor         & demand-side actors
                         (e.g. customer, consumer, competitor, retailer) \\
ProductMaterial     & physical or digital goods
                         (e.g. manufactured product, raw material, 
                          intermediate good, software)\\
SupplyChainEntity   & abstract supertype; used in relation constraints
                         only when no more specific type applies.\\
\end{tabular}
\end{table}

\begin{table}[h]
\centering 
\caption{Relation Types}
\label{tab:relation_types}
\begin{tabular}{lll}
\textbf{Relation Type} & \textbf{Domain} & \textbf{Range}\\
is owned by               & CoreEnterprise      & OrganizationalUnit\\
  has subsidiary            & CoreEnterprise      & OrganizationalUnit\\
  is in partnership with    & CoreEnterprise      & OrganizationalUnit\\
  is regulated by           & CoreEnterprise      & OrganizationalUnit\\
  is an investor in         & CoreEnterprise      & CoreEnterprise\\
  is the owner of           & CoreEnterprise      & ProductMaterial\\
  is distributed by         & ProductMaterial     & SupplyChainPartner\\
  was commissioned by       & ProductMaterial     & CoreEnterprise\\
  typically sells           & CoreEnterprise      & ProductMaterial\\
  sources energy from       & ProductMaterial     & SupplyChainEntity\\
  is manufactured by        & ProductMaterial     & CoreEnterprise\\
  is made from              & ProductMaterial     & ProductMaterial\\
  is designed by            & ProductMaterial     & CoreEnterprise\\
  processes raw material    & SupplyChainEntity   & ProductMaterial\\
  is maintained by          & ProductMaterial     & CoreEnterprise\\
\end{tabular}
\end{table}

\section{SC-RED Dataset Sample}
\label{app:scred_sample}

The following sample has been taken from the training data to represent the format of each sample in the SC-RED dataset. 

\begin{verbatim}
{
    "org_id": "Q5323498",
    "org_label": "EMG 85",
    "text": "The EMG 85 is a popular active humbucker guitar pickup manufactured 
    by EMG, Inc. It is  paired with the 81 in the Zakk Wylde signature EMG set. 
    It was originally designed to be used in the bridge position but is typically 
    installed in the neck position by modern guitar producers.
    ...(content removed for brevity)",
    "spo\_list": [
      {
        "subject": "EMG 85",
        "predicate": "is manufactured by",
        "object": "EMG, Inc."
      }
    ]
}
\end{verbatim}

\section{SC-RED Rudimentary Baseline Prompt} 
\label{app:scred_baseline_prompt}

We present the prompt for the baseline in SC-RED experiments. 

\begin{quote}
{
SYSTEM\_PROMPT = 
\itshape
"""
You are a supply chain relation extraction system. Extract relation triples from the input text.


ONTOLOGY SCHEMA

Entity Types:
  CoreEnterprise       — profit-seeking companies
                         (e.g. logistics company, freight forwarder, 
                          public company, defense contractor)
  SupplyChainPartner   — external suppliers and service providers
                         (e.g. vendor, distributor, OEM, 
                          contract manufacturer, warehouse operator)
  OrganizationalUnit   — internal divisions and regulatory bodies
                         (e.g. government agency, regulatory body, 
                          joint venture, consortium)
  MarketActor          — demand-side actors
                         (e.g. customer, consumer, competitor, retailer)
  ProductMaterial      — physical or digital goods
                         (e.g. manufactured product, raw material, 
                          intermediate good, software, hardware component)
  SupplyChainEntity    — abstract supertype; used in relation constraints
                         only when no more specific type applies.
                         Do not use this type to label entities directly.

Relations (relation\_name: subject\_type → object\_type):
  is owned by               : CoreEnterprise      → OrganizationalUnit
  has subsidiary            : CoreEnterprise      → OrganizationalUnit
  is in partnership with    : CoreEnterprise      → OrganizationalUnit
  is regulated by           : CoreEnterprise      → OrganizationalUnit
  is an investor in         : CoreEnterprise      → CoreEnterprise
  is the owner of           : CoreEnterprise      → ProductMaterial
  is distributed by         : ProductMaterial     → SupplyChainPartner
  was commissioned by       : ProductMaterial     → CoreEnterprise
  typically sells           : CoreEnterprise      → ProductMaterial
  sources energy from       : ProductMaterial     → SupplyChainEntity
  is manufactured by        : ProductMaterial     → CoreEnterprise
  is made from              : ProductMaterial     → ProductMaterial
  is designed by            : ProductMaterial     → CoreEnterprise
  processes raw material    : SupplyChainEntity   → ProductMaterial
  is maintained by          : ProductMaterial     → CoreEnterprise


INSTRUCTIONS

Read the input text carefully. Using ONLY the relations defined in the ontology
above, extract all relation triples that are explicitly supported by the text.

- Every predicate must be one of the relations in the ontology.
- Only extract a triple if the text clearly expresses that relation.
- Do not invent new relation names or entity types.
- The subject and object must satisfy the type constraints implied by the ontology.
  (For example, "is manufactured by" must have a ProductMaterial as subject
   and a CoreEnterprise as object.)
- The subject and object strings must be exact text spans from the input.
- Do not infer or hallucinate entities or relations that are not present in the text.
- If no triples are found, return an empty spo\_list.

Follow the format and typing behaviour shown in the example(s) you are given.
Return ONLY the JSON object below, with no additional text.

{
  "spo\_list": [
    {"subject": "...", "predicate": "...", "object": "..."}
  ]
}
"""
}
\end{quote}


\begin{quote}
{
USER\_PROMPT = 
\itshape
"""
Extract all relation triples from the following text:

{chunk\_text}
"""
}
\end{quote}



\section{CRA Prompt Examples}
\label{app:cra_prompt_example}

System and User Prompts for the proposed new architecture called CRA appear below.  

\textbf{CRA Config D - Staged ReAct}

\begin{quote}
{
SYSTEM\_PROMPT = 
\itshape """

You are a supply chain relation extraction agent. Your task is to extract 
relation triples from supply chain text chunks.


Reason and act using the following loop:

  Thought: [brief reasoning about what to do next - 1 to 2 sentences only]
  Action: [the action to take — either RetrieveExamples or Finish]
  ActionInput: [the input to the action]
  Observation: [the result returned to you after an action — do not write this yourself]

ONTOLOGY SCHEMA

Entity Types:
  CoreEnterprise       — profit-seeking companies
                         (e.g. logistics company, freight forwarder, 
                          public company, defense contractor)
  SupplyChainPartner   — external suppliers and service providers
                         (e.g. vendor, distributor, OEM, 
                          contract manufacturer, warehouse operator)
  OrganizationalUnit   — internal divisions and regulatory bodies
                         (e.g. government agency, regulatory body, 
                          joint venture, consortium)
  MarketActor          — demand-side actors
                         (e.g. customer, consumer, competitor, retailer)
  ProductMaterial      — physical or digital goods
                         (e.g. manufactured product, raw material, 
                          intermediate good, software, hardware component)
  SupplyChainEntity    — abstract supertype; used in relation constraints
                         only when no more specific type applies.
                         Do not use this type to label entities directly.

Relations (relation\_name: subject\_type → object\_type):
  is owned by               : CoreEnterprise      → OrganizationalUnit
  has subsidiary            : CoreEnterprise      → OrganizationalUnit
  is in partnership with    : CoreEnterprise      → OrganizationalUnit
  is regulated by           : CoreEnterprise      → OrganizationalUnit
  is an investor in         : CoreEnterprise      → CoreEnterprise
  is the owner of           : CoreEnterprise      → ProductMaterial
  is distributed by         : ProductMaterial     → SupplyChainPartner
  was commissioned by       : ProductMaterial     → CoreEnterprise
  typically sells           : CoreEnterprise      → ProductMaterial
  sources energy from       : ProductMaterial     → SupplyChainEntity
  is manufactured by        : ProductMaterial     → CoreEnterprise
  is made from              : ProductMaterial     → ProductMaterial
  is designed by            : ProductMaterial     → CoreEnterprise
  processes raw material    : SupplyChainEntity   → ProductMaterial
  is maintained by          : ProductMaterial     → CoreEnterprise

YOUR PROCESS

You MUST follow these steps in order:

STEP 1 — RELATION IDENTIFICATION (always first)
  Read the chunk carefully.
  Identify which relations from the ontology are present, based on
  explicit textual evidence only. Do not invent new relations.
  Do not extract triples yet.

  Thought: [reasoning about which relations appear in the chunk]

STEP 2 — RETRIEVAL (optional, only after Step 1)
  If you are unsure how to identify subject/object entities for the
  relations you found, or the chunk involves unfamiliar phrasing,
  call the retriever tool exactly once.

  Thought: I need a few more similarly annotated examples to guide my extraction.
  Action: RetrieveExamples
  ActionInput: [the full input chunk text]
  Observation: [examples will be returned here by the system]

  Do NOT call this if the chunk is clear and straightforward.
  Call at most ONCE. Only AFTER Step 1.
  Do NOT go to Step 3 until you have the observations from this step. 

STEP 3 — TRIPLE EXTRACTION (always last)
  Using ONLY the relations identified in Step 1, extract all triples
  supported by explicit textual evidence.
  - subject and object must be exact entity mentions from the chunk
  - subject and object must satisfy the type constraints in the ontology schema
  - Do not infer or hallucinate entities not present in the chunk
  - Use the reference examples and any retrieved examples as guidance

  Thought: I have identified the relations and have enough context.
           I will now extract the triples and output the final answer.
  Action: Finish
  ActionInput:
  Present relations: [relation\_1, relation\_2, ...]
  {
    "spo\_list": [
      {
        "subject": "...",
        "predicate": "...",
        "object": "..."
      }
    ]
  }


HARD RULES


  - Always complete Step 1 before anything else.
  - Never call RetrieveExamples more than once.
  - Never write an Observation yourself — only the system produces those.
  - The final action must always be Action: Finish.
  - Do not write anything after the ActionInput of Finish.
  - If no triples are found, return an empty spo\_list.

"""
}
\end{quote}

\begin{quote}
{
USER\_PROMPT = 
\itshape 
"""

Here is a real-world example as reference: 

{retrieved\_examples}

Now please begin the task on the following input\_text: 

{chunk\_text}

"""
}
\end{quote}

\textbf{CRA Config C - Flat FSL with dynamic retrieval}

We only share the system prompt for Config C since user prompt is exactly the same as seen in Config D. 

\begin{quote}
{
SYSTEM\_PROMPT\_C = 
\itshape
"""
You are a supply chain relation extraction system. Extract relation triples from the input text.


ONTOLOGY SCHEMA

Entity Types:
  CoreEnterprise       — profit-seeking companies
                         (e.g. logistics company, freight forwarder, 
                          public company, defense contractor)
  SupplyChainPartner   — external suppliers and service providers
                         (e.g. vendor, distributor, OEM, 
                          contract manufacturer, warehouse operator)
  OrganizationalUnit   — internal divisions and regulatory bodies
                         (e.g. government agency, regulatory body, 
                          joint venture, consortium)
  MarketActor          — demand-side actors
                         (e.g. customer, consumer, competitor, retailer)
  ProductMaterial      — physical or digital goods
                         (e.g. manufactured product, raw material, 
                          intermediate good, software, hardware component)
  SupplyChainEntity    — abstract supertype; used in relation constraints
                         only when no more specific type applies.
                         Do not use this type to label entities directly.

Relations (relation\_name: subject\_type → object\_type):
  is owned by               : CoreEnterprise      → OrganizationalUnit
  has subsidiary            : CoreEnterprise      → OrganizationalUnit
  is in partnership with    : CoreEnterprise      → OrganizationalUnit
  is regulated by           : CoreEnterprise      → OrganizationalUnit
  is an investor in         : CoreEnterprise      → CoreEnterprise
  is the owner of           : CoreEnterprise      → ProductMaterial
  is distributed by         : ProductMaterial     → SupplyChainPartner
  was commissioned by       : ProductMaterial     → CoreEnterprise
  typically sells           : CoreEnterprise      → ProductMaterial
  sources energy from       : ProductMaterial     → SupplyChainEntity
  is manufactured by        : ProductMaterial     → CoreEnterprise
  is made from              : ProductMaterial     → ProductMaterial
  is designed by            : ProductMaterial     → CoreEnterprise
  processes raw material    : SupplyChainEntity   → ProductMaterial
  is maintained by          : ProductMaterial     → CoreEnterprise


INSTRUCTIONS

Read the input text carefully. Using ONLY the relations defined in the ontology
above, extract all relation triples that are explicitly supported by the text.

- Every predicate must be one of the relations in the ontology.
- Only extract a triple if the text clearly expresses that relation.
- Do not invent new relation names or entity types.
- The subject and object must satisfy the type constraints implied by the ontology.
  (For example, "is manufactured by" must have a ProductMaterial as subject
   and a CoreEnterprise as object.)
- The subject and object strings must be exact text spans from the input.
- Do not infer or hallucinate entities or relations that are not present in the text.
- If no triples are found, return an empty spo\_list.

Follow the format and typing behaviour shown in the example(s) you are given.
Return ONLY the JSON object below, with no additional text.

{
  "spo\_list": [
    {"subject": "...", "predicate": "...", "object": "..."}
  ]
}
"""
}
\end{quote}


\textbf{CRA Configs A, B}

These configs only change the user prompt, the system prompt is the same as the one seen in Config C. 

\begin{quote}
{
USER\_PROMPT\_B = 
\itshape
"""
Here is a reference example:

{static\_example}

Now extract triples from the following text:

{chunk\_text}
"""
}
\end{quote}

\begin{quote}
{
USER\_PROMPT\_A = 
\itshape
"""
Extract all relation triples from the following text:

{chunk\_text}
"""
}
\end{quote}



\section{Open Source Model Selection}
\label{app:model_selection}

Table~\ref{tab:model_selection} lists the F1, Recall (R), Precision (P) for each and every open source model that was available in GWDG as of February 2026. Both ReAct\_FSL and ReAct\_Memory were tested. glm-4.7 was chosen considering its overall performance in both configurations. 

\begin{table}[h]
\centering
\caption{Open Source Model performance on DuIE2.0 dataset in ReAct\_FSL and ReAct\_Memory configs.}
\label{tab:model_selection}
\begin{tabular}{lll}\textbf{Open Source Model} & \textbf{ReAct\_FSL {F1, R, P}}& \textbf{ReAct\_Memory {F1, R, P}}\\
qwen3-235b-a22b & {0.3849, 0.451, 0.3358} & {0.2439, 0.1471, 0.7143} \\
llama-3.3-70b-instruct & {0.1629, 0.1765, 0.1513}& {0, 0.0, 0.0} \\
gemma-3-27b-it & {0.3188, 0.3235, 0.3143} & {0.2368, 0.1765, 0.36}  \\
qwen3-32b & {0, 0.0, 0} & {0.3136, 0.3627, 0.2761}\\
apertus-70b-instruct-2509 & {0, 0.0, 0.0} & {0, 0.0, 0}\\
devstral-2-123b-instruct-2512 & {0.3604, 0.3922, 0.3333} & {0.1793, 0.1275, 0.3023}\\
deepseek-r1-distill-llama-70b & {0.117, 0.1078, 0.1279} & {0.1572, 0.1765, 0.1417}\\
\textbf{glm-4.7} & \textbf{{0.3565, 0.402, 0.3203}} & \textbf{{0.5, 0.5098, 0.4906}} \\
internvl3.5-30b-a3b & {0.2357, 0.3039, 0.1925} & {0, 0.0, 0.0}\\
llama-3.1-sauerkrautlm-70b-instruct & {0.3103, 0.3529, 0.2769} & {0.0194, 0.0098, 1.0}\\
medgemma-27b-it & {0.1202, 0.1373, 0.1069} & {0.1271, 0.1471, 0.1119}\\
meta-llama-3.1-8b-instruct & {0.0902, 0.0588, 0.1935} & {0, 0.0, 0}\\
mistral-large-3-675b-instruct-2512 & {0.318, 0.3725, 0.2774} & {0.2027, 0.1471, 0.3261}\\
qwen3-30b-a3b-instruct-2507 & {0.3636, 0.4314, 0.3143} & {0.1224, 0.0882, 0.2}\\
qwen3-30b-a3b-thinking-2507 & {0.3229, 0.3529, 0.2975} & {0.4, 0.4314, 0.3729}\\
qwen3-omni-30b-a3b-instruct & {0.1991, 0.2255, 0.1783} & {0.1091, 0.0588, 0.75}\\
& & \\
\end{tabular}
\end{table}

\section{Error Samples - Baseline Analysis on DuIE2.0}
\label{app:error_analysis_baseline_ablation}

We illustrate the various error modes in baseline ablation using the following
examples. Note that since the examples in DuIE2.0 are in chinese, we used DeepL and Google Translate to translate the examples, the thought loops were however, in english, needing no translation. 

\paragraph{Example 1: Recursive-Verification Trap. Example used AgentRE's output for example 10 in test set in attempt 1}

\textbf{Input} "Zhou Yulou lives in constant fear and suspicion; his life seems to have fallen into a chaotic black hole... Episode 56: Synopsis, Cast \& Crew, Character Profiles, Behind-the-Scenes. The TV series ‘Dear, I'm Sorry’ is the second Chinese TV drama filmed by South Korean actor Joo Jin-mo." \\ 
~\\
\textbf{Actions Taken} GetTaskDescription, RetrieveCorrectMemory. Structure for each of these tool calls is Action: .., ActionInput: ..., Observation: ... \\
~\\
\textbf{Analysis Starts.} Notice how the model keeps looking at the subjects like Joo Jin-mo to verify if the subject makes sense in the relation that it has already established. This analysis could also look like an unending repetition of a few tokens, endless reverification and so on, all classified in this category. 

\begin{verbatim}
1. **Analyzing user request:**
    * **Input sentence:** "Zhou Yulou lives in constant fear and suspicion; his life 
    seems to have fallen into a chaotic black hole ……Episode 56: Cast and Crew, 
    Cast List, Character Profiles, Behind-the-Scenes. The TV series *Dear, I'm Sorry* 
    is the second Chinese TV drama filmed by Korean actor Joo Jin-mo"
    * **Task:** Identify relation triples that conform to the specified format.
    * **Available tools:** `Finish`, `GetTaskDescription`, 
    `GetRelationDefinition`, `RetrieveCorrectMemory`. 
    * **Output Format:** JSON containing “spo\_list”.

2. **Analyze the Input Sentence:**
    * “Zhou Yulou lives every day in fear and suspicion; life seems to have fallen into 
    a chaotic black hole...” -> This appears to be a plot description or character 
    status. It mentions “Zhou Yulou”.
    * "56-episode plot, cast cast list, cast member, actor list, character introduction, 
    behind-the-scenes“ -> These are metadata/navigation links on a webpage and are 
    typically not triples about the work itself, unless ”56" refers to the episode 
    number (but usually “episode number” is not in the list of allowed relationships).
    * “The TV series ‘Dear, I'm Sorry’ is the second Chinese TV series starring Korean 
    actor Joo Jin-mo” -> This is the core information.
        * Subject: TV series ‘Dear, I'm Sorry’
        * Relationship: Starring... (implied ‘starring’ or “directed by” or 
        simply “actor”). Let's look at the allowed relationships.
        * Object: Joo Jin-mo 
        * The sentence says "Joo Jin-mo filmed .“ Typically, if an actor appears in a 
        drama, they are the ‘lead’ or ”actor." Let’s check the list of allowed 
        relationships. 
        * Allowed relationships: `Official language|Dynasty|Chairman|Box office|Composer|
        Singer|Lead role|Area|Capital|Screenwriter|Abbreviation|Nationality|Climate|
        Awards|Number|Theme song|Father|Production company|Ancestral home|Guest|Director|
        Album|Spokesperson|Registered capital| Portrays|Wife|Adapted from|Principal|
        Lead Actor|ZIP Code|Mother|City|Area|Population|Producer|Elevation|Years of Study|
        Founder|Voice Actor|Headquarters|Host|Lyricist|Release Date|Establishment Date|
        Husband|Professional Code|Alma Mater|Author`.
        * “Joo Jin-mo” is a “South Korean actor.” The sentence says “It was filmed by 
        South Korean actor Joo Jin-mo.” This implies he starred in it . The relationship 
        “starring” seems most appropriate, or it could be “director” if “filmed” implies 
        directing, but “actor” usually means performing. Let’s look at “Zhou Yulou.” 
        He might be the ‘lead’ or “play” a certain role. 
        However, the sentence only says “Zhou Yulou lives in...”. It does not 
        explicitly state that “Zhou Yulou is a character in *Dear, I'm Sorry*”. 
        Although the context implies this, strictly speaking, the sentence structure 
        keeps them separate.
        * Let’s focus on the explicit statement: “The TV series *Dear, I'm Sorry* 
        was filmed by South Korean actor Joo Jin-mo.”
        
* Subject: *Dear, I'm Sorry* 
        * Object: Joo Jin-mo
        * Predicate: starred in. Why? Because he is an actor. If he were the director, 
        we would typically say “directed by...”. “Starring” is a bit ambiguous for 
        an actor, but considering “actor,” “starring” is the closest match. 
        Let’s check if “directed” would be more appropriate. 
        Usually, ‘starring’ can imply directing. But “Korean actor” 
        strongly points to acting. Let’s look at examples from memory.
        * Memory Example 1: “*While You Were Sleeping*... Lee Jong-suk... starred.”
        * Memory Example 3: “The CEO... directed by Oh Tae-myung”.
        * The sentence says “It was filmed by Korean actor Joo Jin-mo”. This is 
        somewhat unusual phrasing. It might mean he starred in it. 
        Let’s assume it means “starred in”.
        * Wait, what about “Zhou Yulou”? “Zhou Yulou lives his life in...” 
        If Zhou Yulou is a character in the drama, could the relationship be 
        “Zhou Yulou” -> “played by” -> “character name”? 
        No, “Zhou Yulou” sounds like a character name. 
        Who plays him? The text doesn’t say. 
        Or “Dear, I’m Sorry” -> “main character” -> “Zhou Yulou”? 
        The text doesn’t explicitly link Zhou Yulou to the title; 
        it’s only implied by context. 
        However, in these tasks, if the sentence describes the plot, 
        the character is usually the “main character.”
        * Let’s check the “main character” relationship.
        * Let’s check the “played by” relationship.
        
* Let’s check the “Starring” relationship.
        * Let’s check the “Director” relationship.
        * Let’s check the “Screenwriter” relationship.
        * Let’s check the “Release Date” relationship.
        * Let’s check the “Production Company” relationship.

    * 
Let’s re-evaluate “The TV series *Dear, I’m Sorry* was filmed by Korean actor 
Joo Jin-mo.”
        * Subject: *Dear, I’m Sorry*
        * Object: Joo Jin-mo
        * Predicate: starred in. This seems to be the most 
        reasonable interpretation, since he is an actor.

    * Let’s re-evaluate “Zhou Yulou”.
        * Text: “Zhou Yulou lives in constant fear and suspicion...”
        * Context: This looks like a plot summary.
        * Relationship: Is Zhou Yulou the lead actor in *Dear, I’m Sorry*?
        * If I extract `{“subject”: “Dear, I’m Sorry”, ‘predicate’: "starring “, ”object“: ”Joo Jin-mo“}`, this is safe.
        * If I extract `{”subject“: ”Dear, I'm Sorry“, ”predicate“: ”main character“, ‘object’: ”Zhou Yulou“}`, this is risky because the text does not explicitly state that ”Zhou Yulou is a character in *Dear, I'm Sorry*," although it is highly likely.
        * However, looking at the provided memory examples, they typically extract 
        explicitly stated relationships.
        * Let’s check if “Joo Jin-mo” is the “director.” The text says “actor Joo Jin-mo.” 
        So “starring” is better than “director.”

    * Is there any other information?
        * “56 episodes” -> number of episodes? Not in the list.
        * “South Korea” -> Nationality? Subject: Joo Jin-mo, Predicate: Nationality, 
        Object: South Korea. The text says “South Korean actor Joo Jin-mo.” 
        This is a valid triplet.
        * Subject: Joo Jin-mo
        * Predicate: Nationality
        
* Object: South Korea

    * So, potential triples:
        1. Subject: *Dear, I'm Sorry*, Predicate: Starring, Object: Joo Jin-mo
        2. Subject: Joo Jin-mo, Predicate: Nationality, Object: South Korea

    * Let's check if “Zhou Yulou” can be extracted.
        * If I assume Zhou Yulou is a character in the drama, I can extract “lead actor”.
        * Let’s look at Memory Example 2: “*Scarlet Heart* is adapted from... 
        Tong Hua...”.
        * Let’s look at Memory Example 5: "*My Father Is a Bench*... Wang Baoqiang. 
        ..starring“.
        * The sentence structure ”Zhou Yulou... [plot] ... the TV series 
        ‘Dear, I'm Sorry’...“ implies a connection. However, strictly speaking, 
        the ”Zhou Yulou" part might be a remnant of the previous sentence 
        or a general description. But generally, in these extraction tasks, 
        if a character’s name is mentioned, they are the “main character.”
        * Let’s check “GetRelationDefinition” to ensure the definitions of 
        “starring” and “nationality” align with my intuition.

3. **Execution Tools:**
    * I need to check the definitions of ‘starring’ and “nationality.”
    * I should also check “Protagonist”.

    * Action: GetRelationDefinition
    * ActionInput: LeadActor|Nationality|Protagonist

4. **Simulated Tool Output:**
    * “LeadActor”: The main actor in a TV show or movie.
    * ‘Nationality’: The country a person belongs to.
    * “Protagonist”
 
: The main character in a film or TV series.

5. **Refine Triples:**
    * Triple 1:
        * Subject: Joo Jin-mo
        * Predicate: Nationality
        * Object: South Korea
        * Evidence: “South Korean actor Joo Jin-mo”
    * Triple 2:
        * Subject: *Dear, I'm Sorry*
        * Predicate: Lead Actor
        * Object: Joo Jin-mo
        * Evidence: “The TV series *Dear, I'm Sorry* was filmed by South Korean 
        actor Joo Jin-mo” (actor filming -> starring).
    * Triplet 3:
        * Subject: *Dear, I'm Sorry*
        
* Predicate: Lead Character
        * Object: Zhou Yulou
        * Evidence: Zhou Yulou is mentioned in a sentence describing the drama. 
        Although the connection is syntactically loose, semantically, 
        he is the lead character of the drama.
        * *Self-correction*: In many strict datasets, if a link is not explicit 
        (e.g., “Zhou Yulou is the lead character of *Dear, I'm Sorry*”), 
        they might skip it. However, " “Zhou Yulou” is a specific name appearing 
        in the description. Let’s look at the output of “RetrieveCorrectMemory.” 
        It does not display similar examples of loose links. Example 1 has 
        “Ding Zaichan (played by Lee Jong-suk),” which is explicit. 
        Example 5 has “Wang Baoqiang...starring.” 
        * Let’s stick with the most explicit ones. 
        “Joo Jin-mo” -> “Nationality” -> “South Korea” is very clear. 
        “Dear, I'm Sorry” -> ‘Lead’ -> “Joo Jin-mo” is clear (actor in the film) .
        * “Zhou Yulou” -> “Lead Role” -> “Dear, I'm Sorry” is inferred. 
        I will include it because it describes the “lead role,” but I will prioritize 
        the explicit ones. In fact, looking at the prompt “Identify triples that meet 
        the specified format requirements from the given sentences” typically means 
        extracting all valid information. Zhou Yulou is the described character.
        * Let’s check if “Zhou Yulou” is the subject of “plays”? No, there is no mention 
        of an actor.
        * Let’s check if “Zhou Yulou” is the object of “main character”? Yes.
        * Let’s check if “Dear, I’m Sorry” is the subject of “main character”? Yes.

    * Let’s check the relationship regarding “Joo Jin-mo” again.
        * Does “filmed” mean ‘directed’? If the text said " filmed by Joo Jin-mo,“ that 
        might mean directing. But it says ”filmed by Korean actor Joo Jin-mo.“ 
        When an actor ”films“ a drama, it usually means starring in it. 
        So ‘starring’ is correct.

    * Let’s check ”Zhou Yulou" again .
        * Sentence: “Zhou Yulou lives in constant fear and suspicion... the TV drama 
        *Dear, I'm Sorry*...”
        * If I had to be precise, I might skip Zhou Yulou, because the 
        text doesn’t explicitly link the show title to the character’s name 
        (e.g., “In *Dear, I’m Sorry*, Zhou Yulou...”). It just places them in the 
        same sentence.
        * However, these sentences come from web scraping, where the first part 
        is a plot summary.
        * Let’s look at the results for “RetrieveCorrectMemory.” They are very clear.
        * I will stick with the clear ones.
        * 1. Joo Jin-moo - Nationality - South Korea
        * 2. *Dear, I’m Sorry*
 
- Starring - Joo Jin-mo

    * Wait, could “Zhou Yulou” be the “author” of “Dear, I'm Sorry”? 
    No, “Living in Fear” is a character description.
    * Could “Zhou Yulou” be the ‘director’ of “Dear, I'm Sorry”? No.
    
* I will stick to clarifying these points.

    * Let’s check the relationship with “Joo Jin-mo” again.
        * Does “filmed” mean “directed”? If the text said “filmed by Joo Jin-mo,” then 
        he might be the director. But it says “It was filmed by Korean actor Joo Jin-mo.” 
        When an actor “films” a drama, it usually means they star in it. 
        So ‘starring’ is correct.

    * Let’s check “Zhou Yulou” again.
        * Sentence: "Zhou Yulou lives in fear and suspicion all day long... the TV drama 
        *Dear, I’m Sorry*... "
        * If I had to be precise, I might skip Zhou Yulou because the text doesn’t 
        explicitly link the show title to the character’s name 
        (e.g., “In *Dear, I’m Sorry*, Zhou Yulou... “). 
        It simply places them in the same sentence.
        * However, these sentences usually come from web scraping, where the 
        first part is a plot summary.
        * Let’s look at the results from ”RetrieveCorrectMemory." They are very clear.
        * I will stick with the clear ones.
        * 1. Joo Jin-moo - Nationality - South Korea
        * 2. *Dear, I’m Sorry* - Starring - Joo Jin-mo

    * Wait, could “Zhou Yulou” be the “writer” of “Dear, I'm Sorry”? No, “Living in Fear” 
    is a character description.
    * Could “Zhou Yulou” be the ‘director’ of “Dear, I'm Sorry”? No.
    * I will stick with the clear ones.
    
* Let’s check if “Joo Jin-mo” is the “director.”
        * If I query “GetRelationDefinition” for “director,” it says “the core 
        creator responsible for bringing the script to the screen in a 
        film or TV series.”
        * The text says “actor Joo Jin-mo filmed.” This is ambiguous. 
        But “actor” is a modifier. Usually, if he were the director, the text 
        would say “Director Joo Jin-moo.” So ‘starring’ is safer.

    * Let’s check if “Joo Jin-moo” as “screenwriter.” No.

    * Let’s check if “Dear, I’m Sorry” has a “release date.” No.

    * Let’s check if “Dear, I’m Sorry” has a “production company.” No.

    * So, I have two strong candidates.
    
* Candidate 1: {“subject”: “Joo Jin-moo”, “predicate”: “nationality”, 
    “object”: “South Korea”}
    * Candidate 2: {‘subject’: “Dear, I'm Sorry”, “predicate”: “stars”,
    “object”: “Joo Jin-mo”}

    * Let’s consider “Zhou Yulou”. If I include it, that would be 
    `{‘subject’: “Dear, I’m Sorry”, “predicate” : “lead”, ‘object’: “Zhou Yulou”}`. 
    * Let’s try to see if there is any ambiguity. “Zhou Yulou” is a character name. 
    The sentence describes his status. Then the drama title is mentioned. 
    In NLP extraction tasks, this proximity usually implies a relationship. 
    However, there is no explicit verb (e.g., “In the drama, Zhou Yulou...”), 
    which makes it somewhat loose.
    * Let’s look at an example of "RetrieveCorrectMemory“ Example 1: ”...The prosecutor, i
    n a nutshell, is ‘the one who gave up the promotion.’ This prosecutor who 
    gave up his promotion is Ding Zaican (played by Lee Jong-suk).“ 
    Here, ”Ding Zaican“ is explicitly linked to ‘prosecutor’ (the role) and the actor.
    * My sentence: ”Zhou Yulou spends his days living in... the TV drama 
    *Dear, I'm Sorry*..."
    
* I will exclude Zhou Yulou to ensure high accuracy/strict adherence, 
unless I find evidence that “lead actor” permits such a loose link. 
The prompt states: “Identify triples of relationships from the given sentences that 
meet the specified format requirements.” 
    * Let’s assume strict extraction.
    
* Wait, is “Joo Jin-mo” the “lead actor”? The text says “It was... filmed by Joo Jin-mo.” 
    * Is “filmed” a synonym for “lead actor”? Not exactly. “Filmed” 
    usually means ‘filmed’ (a movie/photo). If an actor " “filmed” a drama, 
    they starred in it.
    * Let’s check if ‘director’ is a better fit. “Filmed by...” usually implies a 
    director. “Filmed by...” implies that the subject (the drama) was filmed by 
    the object (the actor). This is odd. It might mean “Joo Jin-mo filmed it.”
    * Let’s check if “Joo Jin-mo” is a “director.” If he’s an actor, maybe he’s also 
    a director? Or perhaps the text is simply poorly worded.
    * Let’s look at the output of “GetRelationDefinition” (simulation).
    * “Lead Actor”: The main actor in a film or TV series.
    * ‘Director’: The core creator responsible for bringing the script to the screen 
    in a film or TV series.
    * Given “actor,” “lead actor” is the most logical mapping.

    * Final decision:
    
* 1. Joo Jin-mo - Nationality - South Korea
    * 2. “Dear, I'm Sorry” - Lead Actor - Joo Jin-mo

    * Let’s check “Zhou Yulou” again. If I include it, it would be 
    `{“subject” : “Dear, I'm Sorry”, “predicate”: “lead role”, ‘object’: “Zhou Yulou”}`. 
    * Let’s try to see if there’s any ambiguity. “Zhou Yulou” is a character name. 
    The sentence describes his status. Then the drama title is mentioned. 
    In NLP extraction tasks, this proximity usually implies a relationship. 
    However, there is no explicit verb (e.g., "in the drama, Zhou Yulou... “), which 
    makes this a bit loose.
    * Let’s look at Example 1 from ”RetrieveCorrectMemory“: ”...The prosecutor, 
    in a nutshell, is ‘the one who gave up the promotion.’ 
    This prosecutor who gave up the promotion is Ding Zaichan 
    (played by Lee Jong-suk).“ 
    Here, ”Ding Zaichan“ is explicitly linked to ‘prosecutor’ (the role) 
    and the actor.
    * My sentence: ”Zhou Yulou spends his days in... the TV drama 
    *Dear, I'm Sorry*...“. 
    * I will exclude Zhou Yulou to ensure high accuracy/strict adherence, unless I 
    find evidence that ”main character“ permits such loose linking. 
    The prompt states: ”Identify triples that meet the specified format 
    requirements from the given sentences."
    * Let’s assume strict extraction.
    * Wait, is “Joo Jin-mo” the “lead actor”? The text says “It was... filmed 
    by Joo Jin-mo.”
    
* Is “filmed” a synonym for “starring”? Not exactly. “Filmed” usually means “filmed” 
(a movie/photo). If an actor “filmed” a drama, they starred in it. 
    * Let’s check if ‘directed’ is better. “Directed by...” usually implies the director. 
    “Filmed by...” implies the subject (the drama) was filmed by the object (the actor). 
    That’s odd. It might mean " Joo Jin-mo filmed it.“ 
    * Let’s check if ”Joo Jin-mo“ is the ”director.“ If he’s an actor, 
    maybe he’s also the director? 
    Or perhaps the text is just poorly worded.
    * Let’s look at ”GetRelationDefinition". 
\end{verbatim}

\paragraph{Example 2: CoT Controllability Failure. Example chosen is: Sample 1 from Attempt 1 of AgentRE\_redux with correct and reflexion memories}

\textbf{Input:} The Youth Theoretical Study Group of the All-China Federation of Industry and Commerce Youth Committee, established in August this year, recently held its first lesson in the Party Building Conference Room of Huaxun Fangzhou Group. The group leader, Wu Guangsheng, Chairman of Huaxun Fangzhou Technology Co., Ltd., the deputy group leader, Li Jingjing, Chairman of Liaoning Qingchuang Entrepreneurship Service Group Co., Ltd., and the deputy group leader, Jiang Nan, Chairman of Jinling Gold Foil Co., Ltd., gathered together to learn from each other.\\
~\\

\textbf{Tool calls: GetTaskDescription, RetrieveCorrectMemory, RetrieveRelevantInfo, RetrieveReflexionMemory}
Structure for each of these tool calls is Action: .., ActionInput: ..., Observation: ...\\
~\\
\textbf{Analysis Starts.} Notice how the thought gets completely interrupted when the format triples step begins

\begin{verbatim}
    1. **Analyze the Request:**
* Input sentence: "The Youth Theoretical Study Group of the All-China Federation of 
Industry and Commerce Youth Committee, established in August this year, recently 
held its first lesson in the Party Building Conference Room of Huaxun Fangzhou 
Group. The group leader, Wu Guangsheng, Chairman of Huaxun Fangzhou Technology 
Co., Ltd., and deputy group leaders, Li Jingjing, Chairman of Liaoning Qingchuang 
Entrepreneurship Service Group Co., Ltd., and Jiang Nan, Chairman of Jinling Gold 
Foil Co., Ltd., met to implement the instructions of the All-China Federation of 
Industry and Commerce leaders and accelerate the work progress.

* Task: Extract relation triples (Subject, Predicate, Object).

* Allowed Predicates: `Official Language|Dynasty|Chairman|Box Office|Composer|
Singer|Leader|Land Area|Capital|Screenwriter|Abbreviation|Nationality|Climate|
Award|Number|Theme Song|Father|Production Company|Ancestral Home|Guest|Director|
Album|Spokesperson|Registered Capital|Played|Wife|Adapted from|Principal|
Leading Role|Postal Code|Mother|City|Area|Population|Producer|Altitude|Years of Study|
Founder|Voiceover|Headquarters Location|Host|Lyricist|Release Date|
Establishment Date|Husband|Professional Code|Graduated Institution|Author`

* Output Format: JSON `{"spo_list": [...]}`.

2. **Analyze the Sentence:**

* "The All-China Federation of Industry and Commerce Youth Committee Youth 
Theoretical Study Group, established in August of this year..." 
-> Entity: "All-China Federation of Industry and Commerce Youth Committee 
Youth Theoretical Study Group", Relation: "Establishment Date", 
Object: "August of this year".

* "...The leader of this theoretical study group, Wu Guangsheng, Chairman of 
Huaxun Fangzhou Technology Co., Ltd..." 
-> Entity: "Huaxun Fangzhou Technology Co., Ltd.", 
Relation: "Chairman", Object: "Wu Guangsheng".

* "...and Deputy Group Leader, Chairman of Liaoning Youth Innovation and 
Entrepreneurship Service Group Co., Ltd., Li Jingjing..." 
-> Entity: "Liaoning Youth Innovation and Entrepreneurship Service Group Co., Ltd.", 
Relation: "Chairman", Object: "Li Jingjing".

* "...and Deputy Group Leader, Chairman of Jinling Gold Foil Co., Ltd., Jiang Nan..." 
-> Entity: "Jinling Gold Foil Co., Ltd.", Relation: "Chairman", 
Object: "Jiang Nan".

3. **Verify Predicates:**

* "Establishment Date" is in the list.

* "Chairman" is in the list.

4. **Construct Triples:**

* (Youth Theoretical Study Group of the All-China Federation of Industry and Commerce 
Youth Committee, Establishment Date, August of this year)

* (Huaxun Fangzhou Technology Co., Ltd., Chairman, Wu Guangsheng)

* (Liaoning Youth Innovation and Entrepreneurship Service Group Co., Ltd., 
Chairman, Li Jingjing)

* (Jinling Gold Foil Co., Ltd., Chairman, Jiang Nan)

5. **Format .
\end{verbatim}

\paragraph{Example 3: Parsing Error. Example chosen is: Sample 20 from Attempt 3 of AgentRE}

We shall be skipping over the rest of the thought process to specifically highlight the parsing error. In the example one can see that the Finish action was called but the model corrected itself which was not parsed by the output parser resulting in a retry.  \\
~\\
\textbf{Input:} The song 'Southward Flying Geese' is included in the album 'New Drunken Concubine'\\
~\\
\begin{verbatim}
    Thought: Based on the task description and recall example, I need to extract 
    relation triples from the sentence. The sentence is "The song 'Southward 
    Flying Geese' is included in the album 'The New Drunken Concubine'". 
    This contains two entities: "Southward Flying Geese" (song) and "The New 
    Drunken Concubine" (album). The relation between them is "Included in", 
    corresponding to "Album" in the relation list. According to the example, 
    entity names do not need to include quotation marks. 
    Therefore, the triple is (Southward Flying Geese, Album, The New Drunken Concubine).

    Action: Finish
    ActionInput: {"spo_list": [{"subject": "Southward Flying Geese", 
    "predicate": "Album", "object": "The New Drunken Concubine"}]}
    </think> This contains two entities: "South Flying Goose" (song) 
    and "New Drunken Concubine" (album). The relationship between them is "included in," 
    corresponding to "album" in the relation list. According to the example, 
    entity names do not need to include quotation marks. 
    Therefore, the triple is (South Flying Goose, album, new Drunken Concubine).
    
    Action: Finish
    ActionInput: {"spo_list": [{"subject": "South Flying Goose", "predicate": "album", 
    "object": "new Drunken Concubine"}]}
    "South Flying Goose" "New Drunken Concubine" (South Flying Goose, album, 
    new Drunken Concubine).
\end{verbatim}

\section{SC-RED Ablation Error Samples}
\label{app:scred_error_samples}

\paragraph{No output after retries}

Since both attempts at solving the task give similar output, we only show one attempt for the example with index 41 in the ablation test set as solved by CRA Config D (Complete Setup) in Run 1. We have removed the bulk of the input text and the triples for the sake of brevity and have only shown this as a representation of the behavior. 

\begin{verbatim}
{
        "idx": 41,
        "system_prompt": "...",
        "user_prompt": "\n\nHere is a real-world example as reference: \n\n\n\n\n\n
        Example 1: \nText: Google AdSense.... Triples: [{'object': 'Google', 
        'predicate': 'is owned by', 'subject': 'Google AdSense'}]\n\n
        ------------------------------------------
        \n\n\n\n
        Now please begin the task on the following input_text: 
        \n\n\"Webtoons (Korean: ; RR: Weptun) are a type of 
        episodic digital comics that originated 
        in South Korea usually meant to be read on smartphones.....",
        "llm_output": "Thought: I need to read the text and identify which relations 
        from  the ontology are present based on explicit textual evidence. 
        I will look for relationships between companies, products, and organizational 
        units mentioned in the text.\n\nAction: Finish\
        nActionInput:\nPresent relations: [is owned by, has subsidiary, 
        is in partnership with, is an investor in, is the owner of, 
        is distributed by, was commissioned by, typically sells, 
        sources energy from, is manufactured by, is made from, 
        is designed by, processes raw material, is maintained by]\n
        {\n  \"spo_list\": [\n    {\n      \"subject\": \"Naver 
        subsidiary Line\",\n      \"predicate\": \"is owned by\",\n      
        \"object\": \"Naver\"\n    },\n    
        {\n      \"subject\": \"Kross Komics\",\n      
        \"predicate\": \"is owned by\",\n      \"object\": \"Kakao Entertainment\"\n},\n    
        {\n      \"subject\": \"Tapas\",\n      \"predicate\": \"is owned by\",
        \n      \"object\": \"Kakao Entertainment\"\n},\n
        {\n      \"subject\": \"Radish Media\",\n      \"predicate\": \"is owned by\"
        ,\n      \"object\": \"Kakao Entertainment\"\n    },\n ......    
        {\n      \"subject\": \"webtoons\",\n      \"predicate\": \"is distributed by\"
        ,\n      \"object\": \"Dong-a Ilbo\"\n    },\n    
        {\n      \"subject\": \"webtoons\",\n      \"predicate\": \"is distributed by\"
        ,\n      \"object\": \"Korea International Trade Association\"\n    },\n"
}
\end{verbatim}
